\chapter{WIP}
\begin{chapquote}{Aristoteles (Fyzika, 239b30)}
Třetí paradox je … že letící šíp je v klidu, což vyplývá z předpokladu, že čas se skládá z okamžiků … Říká, že pokud je vše, co zabírá pořád stejný prostor, v klidu, a pokud to, co je v pohybu, se vždy nachází v „nyní“, pak je letící šíp nehybný.  
\end{chapquote}

\begin{chapquote}{Michael Atiyah}
	Za jasného denního světla matematici prověřují své rovnice a důkazy a při hledání matematické přesnosti nenechávají nic náhodě. V noci však, za svitu úplňku, sní, vznášejí se mezi hvězdami a žasnou nad zázrakem nebes. Nacházejí inspiraci. Bez snů není umění, není matematika, není život.
\end{chapquote}


Stojí před námi hádanka, která lidstvo trápila 2000 let. Antickým filosofům z Eleje se zdál pojem pohybu tak zvláštní, že ho zavrhli a řekli, že žádný pohyb vlastně neexistuje. Zenón Elejský vymyslel několik paradoxů, které měli prokázat logickou nesmyslnost pohybu. Bohužel se nedochovali jeho původní spisy a musíme čerpat z Aristotelova komentáře. Aristoteles mluví o čtyřech paradoxech, přičemž ten čtvrtý je velmi vágní, pojďme si představit první tři.

\begin{paradox}[První Zenónův paradox: Dichotomie (dělení na poloviny)]\label{paradox:zenon 1}
	Závodíte na Olympijských hrách ve sprintu na \SI{100}{m}. Abyste uběhli \SI{100}{m}, musíte uběhnout prvních \SI{50}{m}. Pak musíte uběhnout dalších \SI{25}{m}, pak dalších \SI{12,5}{m} atd. Po $n$ krocích tohoto procesu uběhnete 
\begin{equation*}
	\SI{50}{m}+\SI{25}{m}+\hdots + \frac{100}{2^{n-1}}\SI{}{m} = \left(100-\frac{100}{2^{n-1}}\right)\SI{}{m} < \SI{100}{m}
\end{equation*}
Tedy po konečně mnoho krocích neuběhnete \SI{100}{m}. Nekonečně mnoho kroků nelze provést, tudíž nelze uběhnout \SI{100}{m}.
\end{paradox}

\begin{paradox}[Druhý Zenónův paradox: Achilles a želva]\label{paradox:zenon 2}
	Achilles běží rychlostí \SI{1}{m/s}, želva \SI{0,1}{m/s}. Závodí spolu, přičemž želva dostane náskok \SI{9}{s}, takže když vystartuje Achilles, želva už urazila \SI{0,9}{m}. Doběhne Achilles želvu? Když Achilles urazí vzdálenost \SI{0,9}{m}, o kterou měla želva náskok, bude mu to trvat \SI{0,9}{s} a želva mezi tím urazí vzdálenost \SI{0,09}{m}. O tuto vzdálenost navíc musí Achilles dobíhat želvu, bohužel, za \SI{0,09}{s} k tomu potřebných želva urazí zase \SI{0,009}{m}. Po $n$ krocích urazí běžci vzdálenosti
\begin{align*}
	\text{Želva:} \quad &\SI{0,9}{m}+\SI{0,09}{m}+\hdots +9\cdot 10^{-n+1}\SI{}{m}+9\cdot 10^{-n}\SI{}{m}\\
	\text{Achilles:} \quad &\SI{0,9}{m}+\SI{0,09}{m}+\hdots +9\cdot 10^{-n+1}\SI{}{m}
\end{align*}
a tudíž Achilles želvu nikdy nedoběhne.
\end{paradox}

\begin{paradox}[Třetí Zenónův paradox: Letící šíp]\label{paradox:zenon 3} Představme si letící šíp. Šíp letí nějaký čas, který se skládá z jednotlivých momentů. V každém jednotlivém momentu šíp neurazí žádnou vzdálenost. Tudíž šíp za celkový čas nemůže urazit žádnou vzdálenost. Pohyb je nemožný. 
\end{paradox}

Všechny tři paradoxy rozdělují pohyb na nekonečně mnoho dílčích částí a pak tvrdí, že na realizaci \uv{pohybu} bychom nějak museli těchto nekonečně mnoho částí složit dohromady, a že to nejde. V paradoxech \ref{paradox:zenon 1} a \ref{paradox:zenon 2} se sčítá nekonečně mnoho čísel, což se později naučíme pomocí pojmu \emph{nekonečné řady}. V paradoxu \ref{paradox:zenon 3} se sčítá \uv{nulový pohyb} přes všechny momenty v čase, a tvrdí se, že se nemůže nasčítat na nenulový pohyb. My se to ale naučíme: této operaci se říká \emph{integrál}.

%TODO: fancy font na vedce/filosofy
\section{Závěr}
\begin{enumerate}
	\item Můžeme provádět teoretickou analýzu řešení:
	\begin{enumerate}[(i)]
	\item Existuje řešení? Je jednoznačné? Je definované globálně nebo jen na konečném časovém intervalu?
	\item Je řešení omezené? Jak se chová pro $t\to \infty?$ (Asymptotika)
	\item Když změníme počáteční podmínky, budou se řešení s různými počátečními podmínkami vzdalovat, oddalovat? Co když změníme koeficienty rovnice? (Stabilita)
	\end{enumerate}
	\item Můžeme napsat počítačový program, který spočítá řešení přibližně. Pak nás zajímá:
	\begin{enumerate}[(i)]
	\item Jak blízko je to, co spočítal program, skutečnému řešení? (Chyba aproximace)
	\item Jak napsat program tak, aby počítal rychle? (Efektivita metody, řád metody)
	\item Pokud má skutečné řešení nějakou vlastnost, například omezenost, má tuto vlastnost i výstup našeho programu? (Numerické metody zachovávající strukturu, $\hdots$)
	\end{enumerate}
\end{enumerate}
Teoretická analýza je úzce provázaná s vývojem numerických metod. Například pokud chci mluvit o tom, že přibližné řešení je \uv{blízko} skutečnému řešení, potřebuji existenci a jednoznačnost skutečného řešení. Pokud rovnice není stabilní, to jest dvě blízká řešení se průběhem času vzdalují, musíme na to dávat pozor v numerickém řešení, protože to pravděpodobně znamená, že s časem bude růst chyba aproximace. Pokud počítám pohyb hmotného bodu, při kterém se zachovává energie, chci použít numerickou metodu, která tento požadavek respektuje.




\begin{definice}[Kuželosečka]
  \item Rovina je množina bodů $(x,y,z)\in \mathbb{R}^{3}$ taková, že $ax+by+cz=d$ pro nějaké $a,b,c,d\in \mathbb{R}.$ Půdorysna je rovina s rovnicí $z=0.$
  \item Kuželová plocha je množina bodů $(x,y,z)\in \mathbb{R}^{3}$ taková, že $z^{2}=k^{2}(x^{2}+y^{2})$ pro nějaké $k>0.$ Povrchové přímky kuželu jsou přímky, které náleží kuželové ploše.
  \item Kuželosečka je průsečík kuželové plochy s rovinou.
\end{definice}

Z definice plyne, že kuželosečka slňuje
\begin{align*}
	cz&=d-ax-by \\
	\Rightarrow  k^{2}\alpha^{2}(x^{2}+y^{2})&=(d-ax-by)^{2}
\end{align*}
a naopak by šlo ukázat, že množina řešení jakékoliv rovnice tvaru
\begin{equation*}
  a_0x^{2}+a_1xy+a_2y^{2}+a_3x+a_4y+a_5=0
\end{equation*}
je buď prázdná, nebo se dá vyjádřit jako kuželosečka. Zatím ale není jasné, co geometricky znamenají koeficienty $a_0,\hdots,a_5.$ Pojďme na tom zapracovat.

\begin{věta}[Dandelin-Queteletova věta]
	Mějme kuželovou plochu $\mathcal{P}$ a rovinu $\rho$. Označme $\alpha$ úhel, který svírá $\rho$ s půdorysnou, a $\beta$ úhel, který svírájí povrchové přímky $\mathcal{P}$ s půdorysnou. Prvně předpokládejme $\alpha\neq \beta.$ Sestrojme obě koule $\kappa_1,\kappa_2$ vepsané $\mathcal{P}$ a $\rho$. Body doteku koulí s rovinou $\rho$ nazvěme $F_1, F_2.$ Průnik koulí a roviny jsou dvě kružnice $k_1:=\kappa_1\cap\rho,\ k_2:=\kappa_2\cap\rho$. Zvolme libovolnou povrchovou přímku $p\subset \mathcal{P}$ a označme $P_1:=p\cap k_1,\ P_2:=p\cap k_2,\ 2a:=|P_1P_2|.$ Pak platí:
  \begin{enumerate}[(i)]
	  \item Pokud $\alpha<\beta,$ každý bod $X$ na kuželosečce $\mathcal{P}\cap \rho$ splňuje
\begin{equation*}
	|XF_1|+|XF_2|=2a \qquad (\text{elipsa})
\end{equation*}
	\item Pokud $\alpha>\beta,$ každý bod $X$ na kuželosečce $\mathcal{P}\cap \rho$ splňuje
\begin{equation*}
	||XF_1|-|XF_2||=2a \qquad (\text{hyperbola})
\end{equation*}
\end{enumerate}
Pokud $\alpha=\beta$, můžeme sestrojit jen jednu vepsanou kouli $\kappa$, $F:=\kappa\cap\rho,$ $k:=\kappa\cap \mathcal{P}$. Nechť $\pi$ je rovina určená kružnicí $k$, označme průsečnici rovin $d:=\pi\cap \rho.$ Pak každý bod na kuželosečce $\mathcal{P}\cap \rho$ splňuje
\begin{equation*}
	|XF|=|Xd| \qquad (\text{parabola})
\end{equation*}
\end{věta}
\begin{Todo}
  obrazek
\end{Todo}
\begin{důkaz}
  Použijeme pozorování, že pro kružnici $k$ a bod $X$ vně $k$ můžeme sestrojit dvě tečny ke $k$ procházející bodem $X$, a body doteku tečen s kružnicí $P_1, P_2$ splňují
  \begin{equation}\label{eq:tecny}
    |PP_1|=|PP_2|
  \end{equation}
  Dokazujme bod (i). Pro bod $X$ na kuželosečce označme $p^X\subset \mathcal{P}$ povrchovou přímku, která prochází bodem $X$. Označme $P_1^X:=p^X\cap k_1,\ P_2^X:=p^X\cap k_2.$ Definujme $2a:=|P_1^XP_2^X|.$ Z \ref{eq:tecny} plyne
  \begin{equation*}
    |XF_1|=|XP_1^X|,\ |XF_2|=|XP_2^X|  
\end{equation*}
  Na druhou stranu z rotační symetrie $\mathcal{P}$ a $k_1,k_2$ plyne 
  \begin{equation*}
    \forall X\in \mathcal{P}\cap \rho\colon\ |P_1^XP_2^X|=|P_1P_2|=2a,
  \end{equation*}
  takže
  \begin{equation*}
    |XF_1|+|XF_2|=|XP_1^X|+|XP_2^X|=|P_1^XP_2^X|=2a.
  \end{equation*}
  Bod (ii) se ukáže stejně, akorát $|P_1^XP_2^X|=||XP_1^X|-|XP_2^X||.$
  Ukažme (iii). 
\end{důkaz}
